\section{Implementation}

\subsection{Platform}

The client is a smartphone browser. SmolVLM-256M-Instruct runs there in
4-bit form: the vision encoder and the decoder use 4-bit weights, on WASM
on an iPhone and on WebGPU where the browser provides it. The generated
sentence and its minimum token probability are available before any image
leaves the phone.

Frames that fall below $\tau_{\mathrm{u}}$ are JPEG-encoded and sent over
Wi-Fi to a laptop. The laptop server is FastAPI. It holds the activity
sentence bank and runs CLIP as the checker. Qwen2-VL is the stronger
generative checker in the real-world comparison. Before the JPEG is
transmitted, the phone sends the planned JPEG size and the confirmation
deadline. The laptop reads its CPU utilization and the previous checking
time, predicts $\hat T_i^{\mathrm{edge}}$, and places the request in the
upload scheduler. The scheduler returns a grant. The phone then uploads the
JPEG, the laptop checks it, and the accept-or-rewrite result comes back.
This path does not change the Wi-Fi MAC or the access point.

\subsection{What is Already Running}

The current prototype already runs this client--server path in the browser.
SmolVLM writes the action phrase on the phone, the page shows that phrase
and the minimum token probability, and the JPEG is uploaded only when the
probability is below the threshold. CLIP on the laptop maps both the phone
sentence and, when present, the picture into the sentence bank. The caption
is replaced only when those two matches differ. A confident phrase sends
the sentence and leaves the picture on the phone.

The predictor, the dynamic uplink budget, and the multi-client scheduler
are not in this prototype yet. The next step is to timestamp the Wi-Fi
transfer and to record checker latency under different CPU loads, JPEG
sizes, and checker models. Those traces train the predictor and drive the
scheduler simulation.

\subsection{Evaluation}

The caption path is tested on the phone and the laptop over the same Wi-Fi
network. One person performs the activities in the sentence bank. For this
real-world test we report on-device generation latency, the fraction of
frames that upload a JPEG, the fraction of captions that are accepted or
rewritten, and a human check of the final caption on a small labeled
subset. The three caption baselines are the phone sentence alone, an upload
on every frame, and the confidence-gated upload.

The scheduler is evaluated in simulation driven by the measured upload and
checking traces, with several phones contending for the laptop. We report
prediction mean absolute error, the fraction of checks that meet the
end-to-end deadline, the 95th percentile confirmation delay, the upload
waiting time, and the same figures as the number of clients grows. The
three schedulers are:
\begin{enumerate}
\item \emph{Round robin}: uploads are served with no deadline and no
checking time;
\item \emph{Static budget}: uplink deadlines use one fixed average checking
time;
\item \emph{Predictive budget}: uplink deadlines use
$\hat T_i^{\mathrm{edge}}$.
\end{enumerate}
The caption comparison and the scheduler comparison are reported
separately.

\subsection{Roles}

Both members integrate the demo, the evaluation, and the report.

Juntao Hu is responsible for the local model, the caption bank, the
acceptance rule, the laptop checker, and the caption-quality evaluation.

Tianchen Liu is responsible for the latency predictor, the dynamic uplink
budget, the upload scheduler, and the multi-client simulation.

The sentence-and-check path is the information-processing part of the
project. Predictive, deadline-aware upload is the networking part.